% MODULE 10: lower strip (three unnumbered sections). Bars are drawn in main.tex (page coordinates, no shift); bars end at y 1162.5, boxes end at y 1304.5.
% MAIN CONTRIBUTIONS x 13-333 | KEY RESULTS x 339-675 | REFERENCES / CONTACT x 681-1148.
% Every statement summarises other panels of this poster (panels 3-9); no new numbers are introduced here.
\begin{scope}
  \hyphenpenalty=10000 \exhyphenpenalty=10000 \tolerance=9999 \emergencystretch=3em
  \newcommand{\TRule}{\\\noalign{\vskip 3pt}\hline\noalign{\vskip 4pt}}
  % ---------------------------------------------------------------- MAIN CONTRIBUTIONS: five items, one line each
  \newcommand{\ContribItem}[4]{% x, y (top), icon name, text
    \fill[PPaleGreen] (#1+10,#2+10) circle (11);
    \csname Icon#3\endcsname{#1+10}{#2+10}{0.46}
    \node[anchor=west,inner sep=0pt,text=PText,align=left,text width=275\PX,font=\FiraMed\fontsize{20.5pt}{22.5pt}\selectfont] at (#1+28,#2+10.4) {#4};}
  \ContribItem{22}{1171}{Warn}{\textbf{Minimal replacement blockers:} the $H_4$ witness is unstable although all its proper subsets are stable.}
  \ContribItem{22}{1198}{Matrix}{\textbf{Exact closure factorization} of the network return into a local and a collective part.}
  \ContribItem{22}{1225}{Wave}{\textbf{Beyond nodal damping:} self-energy and signed damping redistribution.}
  \ContribItem{22}{1252}{Clock}{\textbf{Grid-embedded PLL gain synthesis} and exact delay transport.}
  \ContribItem{22}{1279}{GraphHop}{\textbf{Information radius on the full IEEE-39:} authority grows, safety must still be verified.}

  % ---------------------------------------------------------------- KEY RESULTS: 4-row table (model | finding | evidence); every number is printed in the panel named in the comment
  % source (TX4 row): panel 3 / module3.tex: 15 of 15 proper subsets stable (derived/TX4_PROPER_SUBSET_CLOSURE.csv), alpha_P4 = 0.12700646782830968 (claims/tx4_summary.json)
  % source (delay row): panel 9 / module9.tex: 41.57 ms (TABLE_D01b_delay_crossings.csv), K_p -14 %, 5 of 5 events (TABLE12_nonlinear_events.csv)
  % source (n-hop row): panel 8 / module8.tex: n_min = 0, 1, 1, 2 for m = 1..4 (E1_nmin.csv); 11 of 12 exact assignments (E5_spillover.csv, graph=frozen)
  % source (delta D_H row): panel 5 / module5.tex: 60 cases (RANK_LAW_SUMMARY.json, TABLE_02_FINITE_DELAY_RANK_LAW.csv)
  \node[anchor=north west,inner sep=0pt,text=PText,font=\FiraMed\fontsize{19.5pt}{21.5pt}\selectfont] at (345,1168) {%
    \setlength{\tabcolsep}{0pt}\renewcommand{\arraystretch}{1}\arrayrulecolor{PBorder}%
    \begin{tabular}{@{}>{\raggedright\arraybackslash}p{58\PX}@{\hspace{5\PX}}>{\raggedright\arraybackslash}p{150\PX}@{\hspace{5\PX}}>{\raggedright\arraybackslash}p{108\PX}@{}}
      {\FiraSemi\color{PGreen}model} & {\FiraSemi\color{PGreen}finding} & {\FiraSemi\color{PGreen}evidence}\TRule
      TX4 IEEE-39 & $H_4$ blocker, 15/15 proper subsets stable, $\alpha_\perp(H_4)=+0.127\,\mathrm{s^{-1}}$ & full-spectrum, fixed policy\TRule
      full IEEE-39, exact delay & first margin crossing 41.57\,ms; uniform $K_p\,{-}14\,\%$ restores margin, 5/5 events & argument-principle counts + Julia events\TRule
      $n$-hop, full IEEE-39 & $n_{\min}=0,1,1,2$ for $m=1$--$4$; 11/12 exact assignments leave roots beyond the margin & modal equations + full counts\TRule
      $\delta D_H$ & 60/60 cases with one positive and one negative eigenvalue & Julia validation
    \end{tabular}};

  % ---------------------------------------------------------------- REFERENCES / CONTACT
  % nine references of plans/REFS_PLAN.md in two columns (Fira Medium 17 pt, hanging indent for the label), contact line across the bottom
  \newcommand{\RefEntry}[2]{\hangindent=1.9em\hangafter=1\noindent\makebox[1.9em][l]{[#1]}#2\par}
  \node[anchor=north west,inner sep=0pt,text=PText,font=\FiraMed\fontsize{17pt}{19pt}\selectfont] at (688,1168)
    {\begin{minipage}[t]{222\PX}\raggedright
      \RefEntry{1}{P. Kundur, \emph{Power System Stability and Control}, McGraw-Hill, 1994.}
      \RefEntry{2}{F. Dörfler et al., ``Sparsity-promoting optimal wide-area control of power networks,'' IEEE Trans. Power Syst., vol. 29, no. 5, pp. 2281--2291, 2014.}
      \RefEntry{3}{L. Huang et al., ``Impacts of grid structure on PLL-synchronization stability of converter-integrated power systems,'' arXiv:1903.05489.}
      \RefEntry{4}{W. Michiels, S. Gumussoy, ``Eigenvalue based algorithms and software for the design of fixed-order stabilizing controllers for interconnected systems with time-delays,'' arXiv:2003.05496.}
     \end{minipage}\hspace{9\PX}\begin{minipage}[t]{222\PX}\raggedright
      \RefEntry{5}{U. Marković et al., ``Understanding small-signal stability of low-inertia systems,'' IEEE Trans. Power Syst., vol. 36, no. 5, pp. 3997--4017, 2021.}
      \RefEntry{6}{X. Wang et al., ``Grid-synchronization stability of converter-based resources---an overview,'' IEEE Open J. Ind. Appl., vol. 1, pp. 115--134, 2020.}
      \RefEntry{7}{F. Dörfler, F. Bullo, ``Kron reduction of graphs with applications to electrical networks,'' IEEE Trans. Circuits Syst. I, vol. 60, no. 1, pp. 150--163, 2013.}
      \RefEntry{8}{S. Skogestad, I. Postlethwaite, \emph{Multivariable Feedback Control: Analysis and Design}, 2nd ed., Wiley, 2005.}
      \RefEntry{9}{A. Ortega et al., ``Graph signal processing: overview, challenges, and applications,'' Proc. IEEE, vol. 106, no. 5, pp. 808--828, 2018.}
     \end{minipage}};
  \node[anchor=west,inner sep=0pt,text=PGreen,font=\FiraSemi\fontsize{17pt}{17pt}\selectfont] at (688,1297.5) {Jorge Luis Mayorga Taborda\ \ {\color{PGold}\textbullet}\ \ Universidad de los Andes\ \ {\color{PGold}\textbullet}\ \ jl.mayorga236@uniandes.edu.co};
\end{scope}
